\chapter{Implementations and Open Problems}\label{ch:implementations}

This chapter looks at the subject from the other side: what has been
programmed, how far the programs go, and what remains to be done. It ends
with a list of what this book itself leaves undone, which is not the same
list.

\section{Bronstein's integrator}\label{sec:impl-axiom}

The algorithm of Parts~I and~II was implemented by Bronstein in Scratchpad,
later Axiom, in 1987 and 1988; his paper of 1990 mentions the implementation
as recently completed~\cite[Conclusions]{bronstein1990}.
Appendix~\ref{app:code} lists its constructors against the chapters of this
book. Read as a design, it has three layers.

The lowest is the function field: a general one and a radical one, the
second with the closed-form integral basis of Chapter~\ref{ch:radical}, and
with them integral bases, over the rationals and over finite fields, and
divisors, in a general and a hyperelliptic form. The middle layer is the
integration of an algebraic function on a curve: the Hermite reduction, the
double resultant for the residues, a search for points of finite order that
reduces to finite fields, and a separate route for genus $0$. The top layer
is the integration of algebraic integrands over transcendental functions,
the mixed case of Part~III.

The split of the bottom layer into a general and a special case, made twice,
is the lesson for an implementer. Simple radical extensions, and among them
square roots, are where every step is explicit, and where most integrands
met in practice live: Trager remarks that restricting to unnested radicals
excludes fewer than one in a hundred of the problems in a table of
integrals~\cite[Appendix~A]{trager1984}. The general case is needed for
completeness.
Bronstein's conclusions add what can be reused unchanged from the algebraic
case when the mixed case is built on it: ``computing normal integral bases,
the genus of a curve, and the order of a divisor, making an extension
regular''.

\section{The state of implementations}\label{sec:impl-state}

FriCAS is the maintained descendant of that code. Its own account of what
the integrator does, on a page of the project's wiki last edited in June
2022~\cite{fricas2022}, is exact and unsparing:
\begin{itemize}
\item ``The core integrator should be complete for purely transcendental
  functions.''
\item When the top-level function is algebraic, ``FriCAS can apply Hermite
  reduction. It can find the logarithmic part only in the purely algebraic
  case. Otherwise, if Hermite reduction does not already yield the integral,
  FriCAS cannot solve it.''
\item Even in the purely algebraic case, ``there are restrictions on the
  Trager resultant: in general FriCAS can only handle resultants having a
  single nonlinear factor''.
\item ``The solver for the Risch differential equation can only handle the
  purely algebraic case. The logarithmic derivative problem is unimplemented
  in the algebraic case.''
\end{itemize}
Its summary is that ``a large part of the Risch algorithm dealing with
algebraic extensions remains unimplemented''. So, by this account, Part~II of this book is
implemented there with a restriction on the residue polynomial, and of
Part~III the Hermite reduction and little else. That was the position more
than thirty years after the paper that Part~III follows.

Of other systems we can say only what the sources we hold say.
Davenport's implementation of his own algorithm, in REDUCE, ``deals only
with square roots''~\cite{davenport1981}. Van Hoeij's paper on integral
bases says that Trager's method was the one used by Maple and by
Axiom~\cite{vanhoeij1994}, and his paper on parametrisation announces his
package of algorithms for algebraic curves for a coming version of
Maple~\cite{vanhoeij1997}. Kauers, in 2008, compares a heuristic of his own
with the integrators of Axiom, Maple and Mathematica, and gives this
judgment of the field: ``all of today's computer algebra systems only
contain a partial implementation of the complete
algorithm''~\cite{kauers2008}. An unpublished note by
Schultz~\cite{schultz2016} revisits Trager's algorithm with a partial
implementation in Mathematica, and says in its abstract that the algorithm
as published ``is incomplete in the sense that assumptions are made about the
function to be integrated in relation to the defining equation'', which the
note undertakes to remove; we have read only that abstract. We have no
source for the present state of the commercial systems and say nothing about
it.

The same gap explains the heuristics. The Gr\"obner-basis method of
Section~\ref{sec:log-groebner}, the substitution methods of
Section~\ref{sec:radical-later} and the parallel integration of
Section~\ref{sec:elem-others} are methods that find many integrals without deciding the question, and the
first of them is offered by its author expressly for a situation in which
the complete algorithm is rarely implemented in full.

\section{Open problems}\label{sec:impl-open}

These are the problems that the sources state as open, or that we could find
no source for. They are stated where they arise; we collect them.

\begin{enumerate}
\item \emph{A rational algorithm for the mixed case.} Bronstein's algorithm
  for an algebraic extension of $k(t)$ uses Puiseux expansions at the special
  places. He names three things a rational algorithm would need: a formula
  for the residues at the special places, a way to integrate in $k$ without
  expansions at those places, and a rational algorithm for the Risch
  equation (Section~\ref{sec:elem-status}).
\item \emph{The Risch equation in the mixed case.} The rational algorithm
  of 1991 finds the denominator of a solution and not the solution
  (Section~\ref{sec:rde-open}).
\item \emph{Cancellation over a nonlinear monomial.} Volume~I states that no
  algorithm is known (Section~\ref{sec:rde-open}).
\item \emph{The parametric logarithmic derivative problem and limited
  integration over a curve.} Decidable by Risch's theorem; we found no
  algorithm stated as such (Section~\ref{sec:param-known}).
\item \emph{Integrals depending on a parameter.} For which values of a
  parameter an integral becomes elementary is a question of arithmetic
  geometry, settled in a modified form by Masser and Zannier and not by an
  algorithm (Section~\ref{sec:fo-parameters}).
\item \emph{The whole algorithm rationally.} The lazy Hermite reduction
  performs the first step without an integral basis. Bronstein asks, at the
  end of his report, ``whether the complete algebraic integration algorithm
  can be performed rationally without computing an integral basis''; he
  observes that the residues can be had from a norm in place of the inner
  resultant, and that the remaining difficulty is to represent divisors and
  test them for principality without desingularising the
  curve~\cite[Conclusions]{bronstein1998lazy}.
\end{enumerate}

\section{What this book has not done}\label{sec:impl-gaps}

A reconstruction should say where it is thin. Appendix~\ref{app:standing}
gives the standing of each result; the larger omissions are these.

\begin{itemize}
\item \emph{Cited, not proved.} The theory of Chapter~\ref{ch:function-fields},
  from Stichtenoth. Trager's criterion for an integral basis and his genus
  formula (Chapter~\ref{ch:integral-bases}). The good-reduction theorems of
  Chapter~\ref{ch:finite-order}, on which the decision procedure rests.
  Chebyshev's theorem. The additive decomposition of Chen, Kauers and
  Koutschan.
\item \emph{Described, not given.} The change of model that removes
  residues from branch places (Section~\ref{sec:log-ideals}). The
  correctness of \seealg{AbsoluteFactor}. Trager's reduction for one radical
  (Section~\ref{sec:radical-trager1979}). Van Hoeij's parametrisation
  without a rational point. Hess's algorithm. Bronstein's rational
  algorithm for the Risch equation. The reduction of divisors modulo a prime
  in detail.
\item \emph{Part~III.} Outside the simple radical case, the integration of
  Chapter~\ref{ch:elementary} and the solution of the Risch equation in
  Chapter~\ref{ch:risch-equation} are summaries of Bronstein's paper and
  Risch's report. A reader who needs them must go to those sources, and for
  the paper to the page images: its text layer does not preserve the
  formulas.
\item \emph{Sources read in part.} The post-2005 work cited in
  Sections~\ref{sec:radical-later} and~\ref{sec:elem-others} was read in
  abstract only, as were several of the papers cited for alternatives. The
  text says so at each place.
\item \emph{Complexity.} There is no analysis of cost anywhere in the book.
  Bronstein's remark stands: the degree of the extension of the constants
  needed to write the answer can be the factorial of the degree of the
  curve, so the worst case is at least exponential, and ``[w]orst-case
  measurements are not very meaningful, however, since non-integrability can
  often be proved very early''~\cite[Conclusions]{bronstein1990}.
\end{itemize}
